% Ejemplo del tema UD. Compilar desde plantillas-latex/:  make ejemplos
\documentclass[aspectratio=169,10pt]{beamer}
\usepackage[spanish,es-noshorthands]{babel}
\usetheme{UD}

\begin{document}

\UDportada{Título de la presentación}
          {Subtítulo o asignatura}
          {Juan Felipe Rodríguez Galindo\\{\fontsize{6.4pt}{7.4pt}\selectfont\mdseries\color{UDSilver}20261595004}}
          {Maestría en Ciencias de la Información y las Comunicaciones}
          {2026}

\begin{frame}{Diapositiva con tarjetas}{Las columnas siempre con onlytextwidth}
  \begin{columns}[T,onlytextwidth]
    \begin{column}{0.48\textwidth}
      \UDCard{white}{UDBorder}{\cardtitle{UDBlue}{Primera idea}{\small Texto de la tarjeta.}}
    \end{column}
    \begin{column}{0.48\textwidth}
      \UDCard{UDLightGray}{UDRed!60}{\cardtitle{UDRed}{Segunda idea}{\small Otro texto.}}
    \end{column}
  \end{columns}
  \vspace{0.4cm}
  \begin{center}
  \begin{tikzpicture}[node distance=0.32cm]
    \node[paso] (a) {\textbf{Paso 1}\\datos};
    \node[paso, right=of a] (b) {\textbf{Paso 2}\\modelo};
    \node[paso, right=of b, fill=UDRed!8, draw=UDRed] (c) {\textbf{Paso 3}\\resultado};
    \draw[flecha] (a) -- (b); \draw[flecha] (b) -- (c);
  \end{tikzpicture}
  \end{center}
  \note{Guion del presentador: no se imprime.}
\end{frame}

\UDcierre{Juan Felipe Rodríguez Galindo}
         {Maestría en Ciencias de la Información\\y las Comunicaciones}
         {Asignatura}

\end{document}
